\scanpage{58}
la formule ci-dessus ; \deleted{Il ne semble \unsure{pas|plus|donc} \unsure{pourtant|possible} que} \inserted{A priori}, la cohomologie stratifiante d'un Module stratifié s'explicite comme une hypercohomologie zariskienne (et encore, dans un complexe dont les \deleted{\unsure{op}} \unsure{différentielles|différentielle} \unsure{n'est|ne sont} pas \unsure{linéaire|linéaires} !), \deleted{et} mais l'inverse ne semble \reversemarginpar\margin{\hfill\raisebox{-10\baselineskip}[0pt][0pt]{\rule{.4pt}{10\baselineskip}}}\normalmarginpar pas vrai. On va voir pourtant que, dans un certain sens, il en est bien ainsi. À peu de choses près, on associe à tt Module quasi-cohérent $F$ (\uline{sans stratification}) un Module \deleted{quasi-cohérent} \inserted{(plutôt, un pro-Module pro-quasi-cohérent)} $Q^{\unsure{\circ|0|\cdot}}(\deleted{\unsure{M}}F)$ stratifié, dont la cohomologie stratifiante est la coh. zariskienne \uline{ordinaire} de $X$. De plus, $Q^\circ(F)$ est un foncteur en $F$, non seulement pour des hom linéaires $F\to G$, mais même \deleted{\unsure{pour un|par des}} \inserted{\unsure{(comme|(cohomo|(conoyaux}} faisceau \deleted{\unsure{additif, au sens de|distributif,|additifs, \illegible{2}}} de $\underline O_X$-modules) pour des \deleted{\unsure{hom}} \inserted{$D$} : $F\mathbin{-}G$ qui sont seulement des op. diff. rel/$S$. \deleted{\unsure{Il faut|En fait} \illegible{3}} \inserted{[-- Dans le cas très} particulier où $\operatorname{diag}_{X/S}:X\longrightarrow X\mathbin{\times}_S X$ est une immersion nilpotente, donc \deleted{\unsure{on a|au} $Q^\circ($} stratification /$S$ = donnée de descente à $S$, $Q^\circ(F)=f^*f_*(F)$].

L'intérêt \inserted{de ces diverses \unsure{vues|voies|voisins},} c'est que l'\uline{hypercoh} \inserted{\uline{zariskienne ordinaire}} \uline{d'un complexe d'op. diff peut s'interpréter comme la coh. stratifiante} \uline{du complexe} correspondant de faisceaux stratifiés ... !
% archive sheet number: 57

